% ============================================================================
% MUSTERLOSUNG: Wiederholung Mi casa
% ============================================================================
% Sequenz 3, Block d: Musterlosung zum Wiederholungs-AB
% ============================================================================

\documentclass[11pt, a4paper]{article}

% ============================================================================
% SW-MODUS TOGGLE
% ============================================================================
\newif\ifswmodus
\swmodusfalse

% ============================================================================
% PAKETE
% ============================================================================

\usepackage[utf8]{inputenc}
\usepackage[T1]{fontenc}
\usepackage[ngerman]{babel}
\usepackage{helvet}
\renewcommand{\familydefault}{\sfdefault}

\usepackage[margin=2cm]{geometry}
\usepackage{enumitem}
\usepackage{tabularx}
\usepackage{booktabs}
\usepackage{multirow}
\usepackage{pgffor}

\usepackage{xcolor}
\usepackage{framed}
\usepackage{tcolorbox}
\tcbuselibrary{skins}

\usepackage{fancyhdr}
\usepackage{lastpage}
\usepackage{needspace}

% ============================================================================
% FARBEN
% ============================================================================

\definecolor{spanisch-color}{HTML}{c41e3a}
\definecolor{grau-color}{HTML}{888888}
\definecolor{hellgrau-color}{HTML}{f5f5f5}
\definecolor{orange-color}{HTML}{b35c00}
\definecolor{loesung-color}{HTML}{c62828}

\definecolor{sw-dark}{HTML}{000000}
\definecolor{sw-gray}{HTML}{666666}
\definecolor{sw-white}{HTML}{ffffff}

\ifswmodus
    \colorlet{spanisch}{sw-dark}
    \colorlet{grau}{sw-gray}
    \colorlet{hellgrau}{sw-white}
    \colorlet{orange}{sw-dark}
    \colorlet{loesung}{sw-dark}
\else
    \colorlet{spanisch}{spanisch-color}
    \colorlet{grau}{grau-color}
    \colorlet{hellgrau}{hellgrau-color}
    \colorlet{orange}{orange-color}
    \colorlet{loesung}{loesung-color}
\fi

\newcommand{\fachfarbe}{spanisch}

% ============================================================================
% VARIABLEN
% ============================================================================

\newcommand{\thema}{Wiederholung: Mi casa}
\newcommand{\sequenz}{03}
\newcommand{\block}{d}
\newcommand{\fach}{Spanisch}
\newcommand{\klasse}{7}
\newcommand{\maxpunkte}{20}
\newcommand{\datum}{\today}

% ============================================================================
% KOPF- UND FUSSZEILE
% ============================================================================

\pagestyle{fancy}
\fancyhf{}
\renewcommand{\headrulewidth}{0.4pt}
\renewcommand{\footrulewidth}{0.4pt}

\lhead{\fach{} -- Klasse \klasse}
\chead{\textcolor{loesung}{\textbf{ML-\sequenz\block{} (MUSTERLOSUNG)}}}
\rhead{\datum}

\lfoot{}
\cfoot{}
\rfoot{Seite \thepage\ von \pageref{LastPage}}

% ============================================================================
% BEFEHLE
% ============================================================================

\newcommand{\punktebox}[1]{\hfill\fbox{\small \_/#1}}
\newcommand{\loesung}[1]{\textcolor{loesung}{\textbf{#1}}}
\newcommand{\luecke}[1]{\rule{#1}{0.4pt}}

\newtcolorbox{merkkasten}[1][]{
    colback=hellgrau,
    colframe=\fachfarbe,
    colbacktitle=white,
    coltitle=\fachfarbe,
    boxrule=1pt,
    arc=0pt,
    left=8pt, right=8pt, top=8pt, bottom=8pt,
    fonttitle=\bfseries,
    attach boxed title to top left={yshift=-2mm, xshift=5mm},
    boxed title style={boxrule=0pt, colback=white},
    title=#1
}

\newtcolorbox{vokabelkasten}{
    colback=white,
    colframe=\fachfarbe,
    boxrule=0.5pt,
    arc=0pt,
    left=5pt, right=5pt, top=5pt, bottom=5pt
}

\newenvironment{aufgabe}[1]{%
    \needspace{5\baselineskip}%
    \nopagebreak[4]%
    \vspace{10pt}
    \noindent\textbf{\textcolor{\fachfarbe}{Aufgabe #1}}
    \nopagebreak[4]%
    \vspace{5pt}
    \nopagebreak[4]%
    \begin{list}{}{\leftmargin=0pt}
    \item[]
}{%
    \end{list}
}

\newcommand{\es}[1]{\textcolor{\fachfarbe}{\textbf{#1}}}

% ============================================================================
% DOKUMENT
% ============================================================================

\begin{document}

% ============================================================================
% KOPFBEREICH
% ============================================================================

\begin{center}
    {\Large\bfseries\textcolor{\fachfarbe}{Arbeitsblatt: \thema}}\\[0.3em]
    {\large\textcolor{loesung}{\textbf{--- MUSTERLOSUNG ---}}}
\end{center}

\vspace{5pt}

\noindent
\begin{tabularx}{\textwidth}{|X|X|X|}
\hline
\textbf{Name:} \textcolor{loesung}{Musterschuler/in} &
\textbf{Klasse:} \klasse &
\textbf{Datum:} \datum \\
\hline
\end{tabularx}

\vspace{15pt}

% ============================================================================
% MERKKASTEN: Raume (mit Losungen)
% ============================================================================

\begin{merkkasten}[Merkkasten: Las habitaciones (Die Raume)]
\textbf{Ubertrage die Vokabeln aus dem Lernpfad:}

\vspace{0.5em}

\begin{tabular}{ll}
\es{la cocina} & = \loesung{die Kuche} \\[0.8em]
\es{el salon} & = \loesung{das Wohnzimmer} \\[0.8em]
\es{el dormitorio} & = \loesung{das Schlafzimmer} \\[0.8em]
\es{el bano} & = \loesung{das Badezimmer} \\[0.8em]
\es{el jardin} & = \loesung{der Garten} \\
\end{tabular}
\end{merkkasten}

\vspace{10pt}

% ============================================================================
% AUFGABE 1: Raume zuordnen (mit Losungen)
% ============================================================================

\begin{aufgabe}{1}
Verbinde die spanischen Raume mit der deutschen Ubersetzung. \punktebox{4}
\end{aufgabe}

\begin{center}
\begin{tabular}{rcl}
\es{la cocina} & \loesung{---} & das Schlafzimmer \loesung{$\rightarrow$ el dormitorio} \\[0.8em]
\es{el salon} & \loesung{---} & der Garten \loesung{$\rightarrow$ el jardin} \\[0.8em]
\es{el dormitorio} & \loesung{---} & die Kuche \loesung{$\rightarrow$ la cocina} \\[0.8em]
\es{el jardin} & \loesung{---} & das Wohnzimmer \loesung{$\rightarrow$ el salon} \\
\end{tabular}
\end{center}

\textit{\textcolor{grau}{Bewertung: Je 1 Punkt pro korrekter Zuordnung (4 Punkte)}}

\vspace{10pt}

% ============================================================================
% AUFGABE 2: hay + Mobel (mit Losungen)
% ============================================================================

\begin{aufgabe}{2}
Erganze die Satze mit \es{hay} und dem passenden Mobel. \punktebox{4}
\end{aufgabe}

\begin{vokabelkasten}
\textit{Mobel:} una cama | un sofa | una mesa | un armario
\end{vokabelkasten}

\vspace{0.5em}

\noindent
a) En el dormitorio \loesung{hay} \loesung{una cama}. \\[0.8em]
b) En el salon \loesung{hay} \loesung{un sofa}. \\[0.8em]
c) En la cocina \loesung{hay} \loesung{una mesa}. \\[0.8em]
d) En el dormitorio \loesung{hay} \loesung{un armario}. \\

\textit{\textcolor{grau}{Bewertung: Je 0,5P fur hay + 0,5P fur passendes Mobel (4 Punkte)}}

\vspace{10pt}

% ============================================================================
% AUFGABE 3: Farben + Adjektivanpassung (mit Losungen)
% ============================================================================

\begin{aufgabe}{3}
Passe die Farbadjektive an das Nomen an. \punktebox{6}
\end{aufgabe}

\begin{vokabelkasten}
\textit{Farben:} rojo/roja | azul | verde | blanco/blanca | negro/negra | amarillo/amarilla
\end{vokabelkasten}

\vspace{0.5em}

\noindent
a) La mesa es \loesung{roja} (rot). \\[0.8em]
b) El sofa es \loesung{azul} (blau). \\[0.8em]
c) La silla es \loesung{amarilla} (gelb). \\[0.8em]
d) El armario es \loesung{blanco} (weiß). \\[0.8em]
e) Las camas son \loesung{verdes} (grun). \\[0.8em]
f) Los sofas son \loesung{negros} (schwarz). \\

\textit{\textcolor{grau}{Bewertung: Je 1 Punkt pro korrekt angepasstes Adjektiv (6 Punkte)}}

\vspace{10pt}

% ============================================================================
% AUFGABE 4: Zimmer beschreiben (Beispiellosung)
% ============================================================================

\begin{aufgabe}{4}
Beschreibe dein Zimmer auf Spanisch. Verwende mindestens 4 Satze mit \es{hay} und Farbadjektiven. \punktebox{6}
\end{aufgabe}

\textbf{Beispiellosung:}

\loesung{En mi dormitorio hay una cama grande. La cama es blanca.}

\loesung{Hay un armario marron. Tambien hay una mesa pequena.}

\loesung{La mesa es negra. Hay una silla azul.}

\loesung{En mi dormitorio no hay sofa.}

\vspace{0.5em}
\textit{\textcolor{grau}{Bewertungskriterien:}}
\begin{itemize}
    \item Mindestens 4 Satze vorhanden: 2P
    \item Korrekter Gebrauch von \textit{hay}: 2P
    \item Farbadjektive korrekt angepasst: 1P
    \item Sprachliche Korrektheit (Genus, Artikel): 1P
\end{itemize}

% ============================================================================
% PUNKTEUBERSICHT
% ============================================================================

\vspace{20pt}
\hrule
\vspace{10pt}

\begin{flushright}
\textbf{Punkte:} \loesung{\maxpunkte} / \maxpunkte
\end{flushright}

\end{document}
